\documentclass[11pt, letterpaper]{article}

% --- ESSENTIAL PACKAGES ---
\usepackage[utf8]{inputenc}
\usepackage[T1]{fontenc}
\usepackage{textcomp}           % For upright mu, etc.
\usepackage[margin=1in]{geometry} % Sensible margins

% --- MATH PACKAGES ---
\usepackage{amsmath}
\usepackage{amssymb}
\usepackage{mathtools}          % Extends amsmath
\usepackage{amsthm}             % For theorem environments

% --- UTILITY PACKAGES ---
\usepackage{graphicx}           % For including graphics
\usepackage{booktabs}           % For professional-quality tables
\usepackage{siunitx}            % For typesetting SI units and numbers
\usepackage[noblocks]{authblk}  % For author affiliations

% --- BIBLIOGRAPHY (MODERN APPROACH) ---
\usepackage[backend=biber, style=numeric-comp, sorting=none]{biblatex}
\addbibresource{references.bib} % Your .bib file

% --- HYPERREF & CLEVEREF ---
\usepackage[colorlinks=true, allcolors=blue]{hyperref} % For hyperlinks

% cleveref must be loaded after hyperref and amsthm
\usepackage[capitalise, noabbrev]{cleveref}

% --- THEOREM-LIKE ENVIRONMENTS SETUP ---
\newtheorem{theorem}{Theorem}[section]
\newtheorem{lemma}[theorem]{Lemma}
\newtheorem{proposition}[theorem]{Proposition}
\newtheorem{corollary}[theorem]{Corollary}
\theoremstyle{definition}
\newtheorem{definition}[theorem]{Definition}
\theoremstyle{remark}
\newtheorem{remark}[theorem]{Remark}


% --- METADATA ---
\title{Project Proposals for Software Development}
\author[1]{Julian L. Avila-Martinez}
\author[1]{Camilo A. Huertas}
\author[1]{Sebastian Rodriguez}
\affil[1]{Programa Académico de Física, Facultad de Ciencias Matemáticas y
	Naturales, Universidad Distrital Francisco José de Caldas, Bogotá D.C.,
	Colombia}

\date{\today}

\begin{document}

\maketitle

\section{Introduction}

\section{Introduction}
This proposal outlines the development of a scientific software package designed
for problem-solving in relativistic quantum mechanics and quantum field theory.
We present two potential project architectures that, while sharing a common
theoretical foundation in geometric algebra, address different computational
objectives. A critical evaluation of the scope and technical limitations of each
path is essential for determining the final development strategy.

The two proposed projects are:
\begin{itemize}
	\item \textbf{A Geometric Algebra QED Toolkit:} A symbolic and numerical
		library for performing calculations in Quantum Electrodynamics (QED) and
		Quantum Field Theory (QFT) founded upon the Space-Time Algebra (STA)
		framework.
	\item \textbf{A Numerical Dirac Equation Solver:} A dedicated numerical solver
		for the Dirac equation, featuring a flexible interface for defining various
		interaction potentials and boundary conditions, with the potential for
		extension to curved spacetime.
\end{itemize}

\section{A Geometric Algebra QED Toolkit}

\subsection{Project Goal}
The primary objective of this project is to develop a robust symbolic and
numerical library for performing calculations in Quantum Electrodynamics (QED)
using the mathematical framework of Space-Time Algebra (STA), the Clifford
algebra $\mathcal{C}\ell_{1,3}(\mathbb{R})$. The software will provide a
high-level interface enabling users to define and manipulate fundamental objects
such as spinors, Dirac operators, and interaction terms directly within the
algebraic structure of STA, abstracting away the conventional matrix-based
representations.

\subsection{Core Features}

\subsubsection{STA Core Engine}
The foundation of the toolkit will be a computational engine implementing the
algebra $\mathcal{C}\ell_{1,3}(\mathbb{R})$. This engine will be built upon a
basis $\{\gamma_0, \gamma_1, \gamma_2, \gamma_3\}$ satisfying the
anticommutation relation $\{\gamma_\mu, \gamma_\nu\} = 2\eta_{\mu\nu}$, where
$\eta_{\mu\nu}$ is the Minkowski metric with signature $(+,-,-,-)$. The engine
must support the full 16-dimensional multivector basis, which includes:
\begin{itemize}
	\item \textbf{Grade 0:} scalars (1 dimension)
	\item \textbf{Grade 1:} vectors (4 dimensions)
	\item \textbf{Grade 2:} bivectors (6 dimensions, e.g., $\gamma_0\gamma_1$)
	\item \textbf{Grade 3:} trivectors (4 dimensions, e.g., $\gamma_1\gamma_2\gamma_3$)
	\item \textbf{Grade 4:} the pseudoscalar $I = \gamma_0\gamma_1\gamma_2\gamma_3$
\end{itemize}

The core engine must implement the fundamental algebraic operations:
\begin{itemize}
	\item \textbf{Geometric Product:} $ab$
	\item \textbf{Inner Product:} $a \cdot b = \frac{1}{2}(ab + ba)$
	\item \textbf{Outer Product:} $a \wedge b = \frac{1}{2}(ab - ba)$
\end{itemize}

\subsubsection{Spinor Representation}
In the STA framework, relativistic spinors are not represented as column vectors
in $\mathbb{C}^4$, but as elements of the even subalgebra,
$\mathcal{C}\ell^+_{1,3}(\mathbb{R})$, which is isomorphic to the Pauli algebra
$\mathcal{C}\ell_{3,0}(\mathbb{R})$. A relativistic spinor $\psi$ is constructed
by applying a Lorentz transformation to a spinor in its rest frame.

Consequently, Lorentz transformations are not realized via matrix
multiplication. Instead, they are represented by algebraic conjugation with a
\textbf{rotor}, an element $R \in \mathrm{Spin}(1,3)$ satisfying $R\tilde{R}=1$:
\[
	\psi' = R\psi
\]
Rotors are generated within the algebra via the exponential map of a bivector $B$:
\[
	R = e^{B/2}
\]
This representation unifies boosts and rotations and is computationally more
efficient and geometrically more insightful than the standard matrix approach.

\subsubsection{The Dirac Equation in STA}
The standard covariant Dirac equation, $(i\hbar\gamma^\mu\partial_\mu - mc)\psi
= 0$, is translated into a remarkably compact and elegant real geometric
equation:
\[
	\nabla\Psi I\sigma_3 = m\Psi
\]
Here, $\nabla = \gamma^\mu\partial_\mu$ is the vector derivative, $\Psi$ is an
8-component real spinor (an element of the even subalgebra), and the bivector
$I\sigma_3$ (where $\sigma_3 = \gamma_1\gamma_2$) assumes the role of the
complex unit $i$. The software must allow for the symbolic definition and
algebraic manipulation of this equation, enabling, for instance, the derivation
of plane-wave solutions and the demonstration of their Lorentz covariance.

\subsubsection{Interaction and Scattering Amplitudes}
The toolkit will be equipped to handle calculations for fundamental QED
processes, such as Compton scattering. The interaction term $j \cdot A$, where
$j = \psi\gamma_0\tilde{\psi}$ is the conserved Dirac current and $A$ is the
electromagnetic four-potential, is expressed naturally within the algebra. A key
advantage of this formalism, which the software will leverage, is that the
calculation of scattering amplitudes, which traditionally requires the tedious
evaluation of traces of products of gamma matrices, is replaced by elegant and
direct algebraic manipulation of multivectors.

\section{A Numerical Dirac Equation Solver}

\subsection{Project Goal}
The objective of this project is to develop a high-performance, modular
numerical solver for the time-dependent Dirac equation in 3+1 dimensions. The
software will serve as a computational laboratory for investigating a range of
relativistic quantum phenomena. A key design principle will be a flexible
architecture that allows for the straightforward implementation of arbitrary
external potentials and boundary conditions, enabling the simulation of fermion
dynamics in diverse physical scenarios, from quantum electrodynamics (QED) to
particle physics in curved spacetime.

\subsection{Core Architecture and Numerical Methods}

The solver will be engineered to find solutions $\psi(\mathbf{x}, t)$ to the
time-dependent Dirac equation:
\[
	i\hbar \frac{\partial \psi}{\partial t} = \hat{H}\psi = \left[
	c\boldsymbol{\alpha} \cdot (\hat{\mathbf{p}} - q\mathbf{A}) + q\phi + \beta
m_0 c^2 \right] \psi
\]
where $\boldsymbol{\alpha}$ and $\beta$ are the standard Dirac matrices, $(\phi,
\mathbf{A})$ is the electromagnetic four-potential, and $\psi$ is a
four-component spinor.

The numerical evolution will be implemented using a \textbf{Split-Operator
Fourier Transform method}. This pseudo-spectral technique is particularly
well-suited for the Dirac equation as it naturally handles the stiff spatial
derivative term ($\boldsymbol{\alpha} \cdot \hat{\mathbf{p}}$) with high
accuracy in momentum space while applying the potential term in position space.
This method offers excellent stability and preserves the unitarity of the time
evolution, which is critical for long-term simulations. The core solver will be
built on a modular C++ or Python backend, with a user-friendly interface for
defining the simulation grid, initial wavepacket conditions, and potential
functions.

\subsection{Key Research Applications}

\subsubsection{Probing QED Phenomena}

The solver's primary application will be the simulation of single-particle QED
effects where an external classical field is dominant. The user can define
arbitrary electromagnetic fields by specifying the scalar potential
$\phi(\mathbf{x}, t)$ and vector potential $\mathbf{A}(\mathbf{x}, t)$.
\begin{itemize}
	\item \textbf{Klein Tunneling:} A key benchmark for the solver will be the
		simulation of Klein tunneling. A sharp potential barrier $V_0$ approaching
		or exceeding the rest-mass energy gap ($V_0 \gtrsim 2m_0c^2$) will be
		implemented. The solver will evolve an incident fermion wavepacket and
		compute the transmission and reflection coefficients. This will allow for a
		direct observation of the paradoxical high-transmission regime, a signature
		of spontaneous pair creation at the barrier, which is a fundamental
		prediction of QED.
	\item \textbf{Wavepacket Dynamics:} The software will enable the study of
		relativistic wavepacket evolution in complex electromagnetic fields, such as
		the cyclotron motion in a uniform magnetic field, the acceleration of a
		spinor in an electric field, and scattering from localized potential wells.
\end{itemize}

\subsubsection{Explorations in Curved Spacetime (Gauge Theory Gravity)}
A long-term and highly innovative goal for this project is to extend the solver
to incorporate gravitational interactions as described by Gauge Theory Gravity
(GTG) or the standard semi-classical formalism. This involves solving the Dirac
equation in a curved spacetime background, where the equation takes the form.
\[
	\left( i\gamma^\mu D_\mu - m \right)\psi = 0
\]
Here, the Dirac matrices $\gamma^\mu(x)$ become position-dependent, and the
derivative $D_\mu = \partial_\mu + \frac{i}{4}\omega_{\mu}^{ab}\sigma_{ab}$ is
the covariant derivative, which couples the spinor to the spacetime geometry via
the spin connection $\omega_{\mu}^{ab}$. The solver's modular design will allow
for the implementation of this covariant derivative, with the spacetime metric
$g_{\mu\nu}(x)$ (and its corresponding spin connection) supplied by the user.

This extension would open the door to simulating phenomena at the intersection
of quantum mechanics and general relativity:
\begin{itemize}
	\item \textbf{Gravitationally-Induced Zitterbewegung:} The solver could track
		the trajectory of a localized wavepacket in a gravitational potential. This
		would allow for a numerical investigation of how the Zitterbewegung---the
		rapid oscillatory motion intrinsic to the Dirac particle---is affected by
		spacetime curvature.
	\item \textbf{Fermions near Black Holes:} By implementing a background metric
		for a non-rotating (Schwarzschild) or rotating (Kerr) black hole, the solver
		could simulate the scattering, absorption, and superradiant effects of a
		fermion wavepacket interacting with the event horizon.
	\item \textbf{Foundations of Particle Creation:} As an ultimate exploratory
		goal, the framework could be used to study the foundational aspects of
		particle creation in strong gravitational fields. By analyzing the evolution
		of quantum modes in the presence of an event horizon, the solver could
		provide insights into the mechanisms underlying phenomena such as Hawking
		radiation, bridging the gap between single-particle relativistic quantum
		mechanics and quantum field theory in curved spacetime.
\end{itemize}

\section{Comparative Analysis and Challenges}

\subsection{Implementation Challenges}
Both proposed projects present distinct and significant technical challenges.
\begin{itemize}
	\item \textbf{QED Toolkit (STA):} The primary challenge lies in the
		development of a robust and efficient symbolic algebra engine. Implementing
		the geometric product and other multivector operations in a computationally
		performant manner is non-trivial. Furthermore, translating high-level QED
		concepts, particularly beyond tree-level processes, into the STA formalism
		requires deep theoretical understanding and sophisticated software
		architecture.
	\item \textbf{Numerical Dirac Solver:} The challenges are predominantly
		numerical and computational. Ensuring the stability and accuracy of the
		time-evolution algorithm, especially for scenarios involving sharp potential
		gradients (Klein tunneling) or complex geometries (curved spacetime), is
		paramount. Managing the significant computational resources (memory and CPU
		time) required for 3D+1 simulations is a major engineering hurdle. The
		implementation of the covariant derivative for gravitational coupling is
		mathematically intensive and requires careful validation.
\end{itemize}

\subsection{Comparison of Project Philosophies}
The two proposals, while rooted in the same physical theory, represent
fundamentally different approaches to computational physics.

\begin{itemize}
	\item \textbf{Methodology:} The QED Toolkit is fundamentally an
		\textbf{algebraic and symbolic} project. Its goal is to create a new
		language for expressing physical laws, focusing on structural elegance and
		analytical manipulation. In contrast, the Dirac Solver is a
		\textbf{numerical and simulation-based} project, aimed at calculating the
		dynamical evolution of quantum systems in specific, often analytically
		intractable, scenarios.
	\item \textbf{Scope and Innovation:} The innovation of the STA toolkit lies in
		its novel \textbf{computational framework}, offering a new paradigm for
		theoretical calculations. Its scope is broad and foundational. The Dirac
		Solver's innovation lies in its \textbf{application to frontier problems},
		pushing the boundaries of what can be simulated. Its scope is deep and
		application-focused.
	\item \textbf{Outcome:} A successful QED Toolkit would yield a powerful new
		tool for theoretical physicists, potentially simplifying complex analytical
		derivations. A successful Dirac Solver would produce a virtual laboratory
		for exploring concrete physical phenomena, yielding data, visualizations,
		and direct insights into particle behavior in extreme conditions.
\end{itemize}

\end{document}
